% Chapter — Camera hook: pose, intrinsics, and landmark projection
% (FR-23). Follows the mandatory FR-29 six-section template: purpose;
% assumptions and domain of validity; derivation; discretization and
% implementation; validation evidence; references. Written ahead of the
% implementation (Phase 6): this chapter is the binding specification for
% the camera-hook module. The implementing change adds the
% chapter-manifest entry and the \testid{} evidence table. The equation
% labels eq:camera:pose, eq:camera:K, eq:camera:landmark,
% eq:camera:apparent, eq:camera:proj, and eq:camera:occlusion are intended
% to be echoed verbatim in the module source (FR-29 traceability).
\chapter{Camera Hook: Pose, Intrinsics, and Landmark Projection}
\label{ch:camera}

\section{Purpose}
\label{sec:camera:purpose}

This chapter specifies the FR-23 camera \emph{hook}: a sensor that emits
geometric truth for offline external rendering and optical-navigation
research --- the camera pose, the pinhole intrinsics, the list of
visible bodies, the apparent Sun vector, and optional landmark pixel
projections --- with \emph{no} in-core rendering and \emph{no} noise. It
is the substrate for the world-model and AI/ML navigation work of the
later phases: an external renderer or a learned perception model
consumes the emitted pose stream and produces images or measurements,
while the core remains small and deterministic. Everything emitted is
truth: Phase~6 exit criterion~7 gates the pose channels \emph{bit-exactly}
against the \texttt{truth} channels, and gates the landmark projections
against an independent NumPy pinhole script to $< 10^{-6}$ pixels ---
which is why every frame, axis, and pixel convention in this chapter is
stated exactly. The criterion also names the intrinsics, which the
\texttt{sensors.camera} record layout does not carry; that clause is
uncovered, and Section~\ref{sec:camera:validation} records why. Direction
outputs carry the FR-23 velocity-aberration correction of
Chapter~\ref{ch:sensors-optical}, applied through the same shared code
path.

\section{Assumptions and domain of validity}
\label{sec:camera:domain}

Assumptions:
\begin{enumerate}
  \item \textbf{Ideal pinhole.} No lens distortion, defocus, rolling
    shutter, or exposure model: projection is the exact central
    projection of equation~\eqref{eq:camera:proj}
    (Hartley and Zisserman~\cite{hartley2004}). Distortion belongs to
    the external renderer or perception model, not the truth hook.
  \item \textbf{Rigid mount, CG-relative.} The camera is rigidly mounted
    at the configured body-frame offset $\vv{r}_{\mathrm{cam}}^{B}$
    \emph{relative to the composite center of mass} (the point whose
    state the \texttt{truth} channels carry), with fixed orientation
    $\quat{q}_{b2c}$. Center-of-mass travel during propellant depletion
    (Chapter~\ref{ch:massprops}) is therefore not re-projected onto the
    mount: for a station-fixed physical camera the neglected offset is
    bounded by the vehicle's CG travel (centimetres to a metre on the
    DX-3 starter fleet) --- immaterial at orbital target ranges, and
    zero for the rigid post-burn spacecraft this hook targets.
  \item \textbf{Spherical occluders.} Visibility against a body uses
    its mean spherical radius (IAU values per Archinal et
    al.~\cite{archinal2018}, as configured); terrain and oblateness are
    not modeled in the occlusion tests. Near-limb landmarks within the
    oblateness band of the true figure may be misclassified; scenarios
    that care sit the landmark away from the limb.
  \item \textbf{Light-time policy.} Per FR-23 and
    Chapter~\ref{ch:sensors-optical}: geometric ephemeris positions,
    velocity aberration on emitted \emph{directions} (landmark
    projections and the Sun vector), residual light-time effects
    bounded there. The emitted \emph{pose} is geometric truth and is
    not aberrated --- aberration is a property of arriving light, not
    of the platform state.
  \item \textbf{Landmarks are body-fixed points.} A landmark is a
    configured position $\vv{\ell}^{T}$ in a target body's body-fixed
    frame (Moon principal-axis, Mars-fixed, or Earth-fixed frames of
    Chapter~\ref{ch:frames}); it rotates rigidly with the body.
\end{enumerate}

Domain of validity: any propagatable state; projections are defined for
landmarks strictly in front of the camera ($Z^{C} > 0$, enforced by the
visibility flag rather than by exception); the model is a total,
noise-free, draw-free function inside the loop and never throws.
Configurations with $f_x, f_y \le 0$, non-positive image dimensions, or
a landmark on an unconfigured body are rejected by the FR-15 validator.

\section{Derivation}
\label{sec:camera:derivation}

\subsection{Pose and frame chain}
\label{sec:camera:pose}

The camera frame $C$ follows the standard computer-vision convention
(Hartley and Zisserman~\cite{hartley2004}): $+Z^{C}$ along the boresight
(out of the camera, toward the scene), $+X^{C}$ toward increasing image
column ($u$, right), $+Y^{C}$ toward increasing image row ($v$, down) ---
a right-handed triad. The chain from inertial to camera is
\begin{equation}
  \boxed{\;
  \quat{q}_{i2c} = \quat{q}_{i2b} \quatmul \quat{q}_{b2c},
  \qquad
  \vv{r}_{\mathrm{cam}}^{I} = \vv{r}^{I}
    + {\dcm{I}{B}(\quat{q}_{i2b})}^{\mathsf T}\, \vv{r}_{\mathrm{cam}}^{B},
  \;}
  \label{eq:camera:pose}
\end{equation}
with $\vv{r}^{I}$, $\quat{q}_{i2b}$ the truth state (composition per
equation~\eqref{eq:notation:quatcomp}). The emitted pose channels are
$\vv{r}^{I}$ and $\quat{q}_{i2b}$ themselves, byte-identical to the
\texttt{truth} group (exit criterion~7).

The constants $\quat{q}_{b2c}$, $\vv{r}_{\mathrm{cam}}^{B}$, and the
intrinsics are not \emph{record} channels --- they do not vary sample to
sample, so repeating them per record would be waste --- but they are
\textbf{echoed once in the log header}, at \texttt{gnc.camera}, as of
SRLOG v1.3. A consumer therefore composes
equation~\eqref{eq:camera:pose} and the intrinsic matrix
equation~\eqref{eq:camera:K} from the log alone, without the mission
file. That is what FR-23 needs: the hook exists to feed an \emph{offline
external renderer}, and the pose channels carry the vehicle state rather
than the camera's, so without the echo a renderer could neither place the
camera nor build its projection matrix from the file. The header rule
that floats never appear there is preserved by carrying each double as
its IEEE-754 binary64 bit pattern in sixteen hex digits, an encoding that
is exact and whose writer runs no float formatter;
\texttt{docs/formats/srlog\_v1.md} section~3 is normative for it. Landmark
\emph{positions} remain resolved-configuration data: the count $L$ already
rides in the record dtype, and echoing a fixture-scale landmark table
would add kilobytes to every camera header.

An earlier revision of this chapter asserted this echo while no
implementation existed --- verified against a run at the time, no
intrinsic, extrinsic, or landmark value appeared anywhere in the header.
That gap was the direct cause of exit criterion~7's uncovered intrinsics
clause; the echo described here is the closure, and
\testid{test_camera_intrinsics_echo_reproduces_logged_pixels} gates it.

\subsection{Pinhole intrinsics}
\label{sec:camera:intrinsics}

The intrinsics are the four pinhole constants $(f_x, f_y, c_x, c_y)$, in
pixels, forming
\begin{equation}
  \mat{K} =
  \begin{bmatrix}
    f_x & 0 & c_x \\
    0 & f_y & c_y \\
    0 & 0 & 1
  \end{bmatrix},
  \label{eq:camera:K}
\end{equation}
with image dimensions $W \times H$ pixels (columns $\times$ rows). Pixel
coordinate convention, normative: $(u, v)$ are continuous image
coordinates with the origin at the \emph{center of the top-left pixel},
$u$ increasing rightward along columns, $v$ increasing downward along
rows; a point exactly on the optical axis projects to $(c_x, c_y)$; the
physical sensor spans $u \in [-\tfrac12,\, W - \tfrac12]$, $v \in
[-\tfrac12,\, H - \tfrac12]$. For a square-pixel camera of horizontal
field of view $2\alpha_h$, $f_x = f_y = (W/2)/\tan\alpha_h$; skew is
identically zero and is not a parameter.

\subsection{Landmark projection chain}
\label{sec:camera:projection}

A landmark $\vv{\ell}^{T}$ fixed to target body $T$ has inertial
position
\begin{equation}
  \vv{r}_{\ell}^{I} = \vv{r}_{T}^{I}
    + {\dcm{I}{T}}^{\mathsf T}\, \vv{\ell}^{T},
  \label{eq:camera:landmark}
\end{equation}
with $\vv{r}_{T}^{I}$ the body's geometric ephemeris position
(Chapter~\ref{ch:ephemeris}) and $\dcm{I}{T}$ its body-fixed rotation
(Chapter~\ref{ch:frames}: ITRF chain for Earth, DE440 principal-axis
frame for the Moon, IAU 2015 elements for Mars). The line of sight from
the camera is
\begin{equation}
  \vv{d}^{I} = \vv{r}_{\ell}^{I} - \vv{r}_{\mathrm{cam}}^{I},
  \qquad
  \hat{\vv{u}}^{I} = \frac{\vv{d}^{I}}{\norm{\vv{d}^{I}}},
  \label{eq:camera:los}
\end{equation}
its apparent direction $\hat{\vv{u}}'^{\,I}$ follows from the
velocity-aberration formula of Chapter~\ref{ch:sensors-optical},
equation~\eqref{eq:optical:aberration}, with the same barycentric
$\vv{\beta}$ of equation~\eqref{eq:optical:beta} (one shared code path,
exit criterion~9), and the camera-frame point is the apparent direction
carried at the geometric range:
\begin{equation}
  \boxed{\;
  \vv{p}^{C} =
    \norm{\vv{d}^{I}}\;
    \dcm{B}{C}(\quat{q}_{b2c})\,\dcm{I}{B}(\quat{q}_{i2b})\,
    \hat{\vv{u}}'^{\,I}
    = [X, Y, Z]^{\mathsf T},
  \;}
  \label{eq:camera:apparent}
\end{equation}
where each DCM is the quaternion-to-DCM map of
equation~\eqref{eq:notation:quat2dcm} applied to the corresponding
frame-transformation quaternion. The pixel projection is the exact
central projection
\begin{equation}
  \boxed{\;
  u = f_x\, \frac{X}{Z} + c_x,
  \qquad
  v = f_y\, \frac{Y}{Z} + c_y .
  \;}
  \label{eq:camera:proj}
\end{equation}

\subsection{Visibility conditions}
\label{sec:camera:visibility}

A landmark is emitted with $\mathrm{visible} = \mathrm{true}$ iff all
four tests pass, evaluated in this order:
\begin{enumerate}
  \item \textbf{In front of the camera:} $Z > 0$ in
    equation~\eqref{eq:camera:apparent}.
  \item \textbf{Within the sensor:} $-\tfrac12 \le u \le W - \tfrac12$
    and $-\tfrac12 \le v \le H - \tfrac12$ per the pixel-center
    convention of Section~\ref{sec:camera:intrinsics}.
  \item \textbf{Near side of the target body:}
    \begin{equation}
      \bigl(\vv{r}_{\mathrm{cam}}^{I} - \vv{r}_{\ell}^{I}\bigr) \cdot
      \bigl(\vv{r}_{\ell}^{I} - \vv{r}_{T}^{I}\bigr) > 0 ,
      \label{eq:camera:nearside}
    \end{equation}
    the camera lies above the landmark's local tangent plane. For a
    spherical body and a surface landmark this test is \emph{exactly}
    equivalent to non-occlusion by the body itself, so no separate
    self-occlusion test exists.
  \item \textbf{Not occluded by any other configured body:} for each
    body with center $\vv{c}^{I}$ and mean radius $R$ (excluding the
    target body, handled by test 3), the segment from camera to
    landmark must miss the sphere:
    \begin{equation}
      s^{*} = \operatorname{clamp}\!\left(
        \frac{(\vv{c}^{I} - \vv{r}_{\mathrm{cam}}^{I}) \cdot \vv{d}^{I}}
             {\vv{d}^{I} \cdot \vv{d}^{I}},\; 0,\; 1\right),
      \qquad
      \text{occluded} \iff
      \bigl\lVert \vv{r}_{\mathrm{cam}}^{I} + s^{*}\vv{d}^{I}
        - \vv{c}^{I} \bigr\rVert^2 < R^2 ,
      \label{eq:camera:occlusion}
    \end{equation}
    the closest point of the (clamped) segment to the sphere center.
    Occlusion geometry is evaluated on the \emph{geometric} positions:
    occlusion is a property of the mass distribution, and at the
    $\lesssim \qty{20.5}{\arcsecond}$ aberration scale the distinction
    only shifts the flag transition by milliseconds of grazing
    geometry.
\end{enumerate}
The \emph{visible-bodies list} applies tests 1, 2, and 4 to each
configured body's center direction (test 3 does not apply to a center);
a body whose center is barely outside the sensor bound while its limb is
inside is listed not-visible --- the list is a coarse scene inventory
for renderers, not a limb model, and the convention is pinned here so
the flag is deterministic.

\subsection{Sun vector}
\label{sec:camera:sun}

The emitted Sun vector is the apparent unit direction to the Sun in
camera axes,
\begin{equation}
  \hat{\vv{u}}_{\odot}^{C}
    = \dcm{B}{C}(\quat{q}_{b2c})\,\dcm{I}{B}(\quat{q}_{i2b})\,
      \hat{\vv{u}}'^{\,I}_{\odot},
  \label{eq:camera:sunvec}
\end{equation}
with $\hat{\vv{u}}'^{\,I}_{\odot}$ aberrated per
equation~\eqref{eq:optical:aberration} --- consistent by construction
with the sun sensor of Chapter~\ref{ch:sensors-optical}, since both
consume the same function. It supports renderer lighting and
phase-angle bookkeeping downstream.

\section{Discretization and implementation}
\label{sec:camera:implementation}

Binding implementation notes for the Phase~6 module:
\begin{enumerate}
  \item \textbf{Module and traceability.} The camera hook is a
    core-side \texttt{ISensor} on the major-cycle schedule (D-5) with
    \emph{no} PRNG stream --- it is noise-free truth and consumes no
    draws. Source comments echo \texttt{eq:camera:pose},
    \texttt{eq:camera:K}, \texttt{eq:camera:landmark},
    \texttt{eq:camera:apparent}, \texttt{eq:camera:proj}, and
    \texttt{eq:camera:occlusion} verbatim.
  \item \textbf{Bit-exact pose emission.} The pose channels are copies
    of the same doubles the \texttt{truth} writer receives --- not
    recomputed values --- so exit criterion~7's \emph{pose} bit-exactness
    clause is satisfied by construction. The intrinsics and the extrinsics
    are taken from the resolved configuration verbatim, without
    recomputation, but they are not written to the record stream: they are
    configuration, and the \texttt{sensors.camera} layout below carries no
    channel for them. The criterion's intrinsics clause is uncovered for
    that reason; see the discussion following the evidence table.
  \item \textbf{Evaluation order.} Per sample: pose copy; per landmark:
    equations~\eqref{eq:camera:landmark}--\eqref{eq:camera:proj}, then
    the four visibility tests in the stated order with short-circuit
    evaluation (the projection values are logged even when a later
    test fails, with the flag false --- consumers filter on the flag);
    then the visible-bodies list; then the Sun vector.
  \item \textbf{Cross-check tolerance.} The $< 10^{-6}$ pixel gate
    against the independent NumPy script absorbs
    double-rounding-order differences between the core and NumPy
    ($\sim 10^{-16}$ relative on kilo-pixel coordinates
    $\approx 10^{-13}$ pixels) with six orders of margin; the gate
    tolerance is set by the reproduction fidelity this chapter's
    conventions must support, not by floating-point noise.
  \item \textbf{Logged channels.} Everything lands under
    \texttt{sensors.cam.*} with unit-suffixed names: pose, per-landmark
    $(u, v, \mathrm{visible})$, per-body visibility, Sun vector.
\end{enumerate}

\section{Validation evidence}
\label{sec:camera:validation}

The module implementing this chapter is
\texttt{cpp/src/sensors/camera.cpp}.

\begin{table}[htbp]
  \centering
  \caption{Validation evidence for the camera hook: core (doctest,
    \texttt{cpp/tests/test\_sensors.cpp}).}
  \label{tab:camera:validation}
  \begin{tabular}{lp{8.2cm}}
    \toprule
    Test ID & Acceptance basis \\
    \midrule
    \testid{sensors_camera_pinhole_projection_and_visibility} &
      Equation~\eqref{eq:camera:proj} on constructed geometry: a landmark
      on the boresight projects to the principal point; a metric camera
      offset moves $u$ by $f_x X/Z$ and $v$ by $f_y Y/Z$, the two shifts
      differing for the same offset under the anisotropic focal lengths
      ($f_x \ne f_y$) exit criterion~7 calls for. Visibility test 1 drops
      the flag when the line of sight runs backward along the boresight;
      test 2 flips it exactly across the half-pixel sensor edge
      $u = W - \tfrac12$ of the pixel-center convention; and test 3,
      equation~\eqref{eq:camera:nearside}, is isolated by one camera pose
      viewing two landmarks on opposite sides of the body center, both in
      front of the camera and both inside the sensor, with only the near
      one flagged visible. \\
    \testid{sensors_factory_kind_vocabulary} &
      The camera kind constructs from the factory at its configured rate
      and reports the \texttt{sensors.camera} group name; non-positive
      intrinsics are refused at construction. \\
    \bottomrule
  \end{tabular}
\end{table}

% longtable rather than table: these rows outgrew a single page once the
% criterion-7 intrinsics clause was gated, and a float too large for the
% page silently drops its trailing rows rather than warning usefully.
\begin{longtable}{lp{8.2cm}}
  \caption{Validation evidence for the camera hook: in-loop behavior and
    the blind pinhole re-gate (pytest,
    \texttt{tests/python/test\_p6\_optical\_gates.py} and
    \texttt{test\_gnc\_missions.py}).}
  \label{tab:camera:validation-py} \\
  \toprule
  Test ID & Acceptance basis \\
  \midrule
  \endfirsthead
  \multicolumn{2}{@{}l}{\itshape Table \thetable\ continued from the previous page.} \\
  \toprule
  Test ID & Acceptance basis \\
  \midrule
  \endhead
  \bottomrule
  \endlastfoot
    \testid{test_full_sensor_suite_reruns_bit_identical} &
      Exit criterion~1 bit identity, camera half. Two seeded runs of the
      reference mission with all six FR-23 sensor kinds declared agree on
      the whole-file SHA-256 and then channel by channel under array bit
      equality, \texttt{sensors.camera} among them. The hook consumes no
      draws, so its contribution is deterministic by construction; the row
      pins that the surrounding sensors' draw schedules do not perturb it. \\
    \testid{test_landmark_pixels_match_independent_reference} &
      Exit criterion~7's pixel clause. Every logged visible-landmark
      $(u, v)$ matches \texttt{tests/refs/pinhole.py} --- a NumPy
      transcription of equations~\eqref{eq:camera:pose}
      and~\eqref{eq:camera:landmark}--\eqref{eq:camera:proj} written from
      this chapter alone, importing nothing from the core --- to
      $< 10^{-6}$ pixels, on a fixture with $f_x \ne f_y$, an off-axis
      principal point, a nonzero mount station, and a non-identity mount
      rotation, so a transposed convention or a shared scale cannot
      hide. The principal point is off-axis only as of the fixture
      revision recorded in the next row. \\
    \testid{test_camera_echo_fixture_is_not_degenerate} &
      Pins the four fixture properties that make each camera constant
      separately resolvable, so the geometry cannot drift into a state
      where one constant stands in for another. The principal point
      previously sat at $(511.5, 383.5)$, which under the pixel-center
      convention is exactly the image center $((W-1)/2, (H-1)/2)$ of the
      $1024 \times 768$ sensor. On that geometry $c_x$ and $c_y$ were
      recoverable from the image dimensions, so discarding the principal
      point entirely and substituting the center changed the worst
      residual by nothing at all --- measured, \num{1.705303e-13} pixels
      both with and without the substitution. The fixture now carries
      $(486, 402)$, against which the same substitution measures
      \num{25.5} pixels. \\
    \testid{test_landmarks_are_actually_visible} &
      Non-vacuity of the row above: the assertion is that more than half of
      all landmark projections pass the reference's own visibility tests,
      so the pixel gate is not satisfied by off-sensor numbers no consumer
      would read. On the committed fixture all \num{300} pass --- the
      criterion-9 off-axis slew of
      chapter~\ref{ch:sensors-optical} happens to hold every landmark
      on-sensor at every sample. \\
    \testid{test_camera_intrinsics_echo_reproduces_logged_pixels} &
      Exit criterion~7's intrinsics clause. The intrinsics and extrinsics
      are taken from the \texttt{gnc.camera} header echo and the pose from
      the \texttt{truth} channels --- both off the log, neither off the
      mission file --- and composed through
      \texttt{tests/refs/pinhole.py} into predicted pixels, required to
      match the logged \texttt{px\_uv} to $< 10^{-6}$ pixels. The row is
      non-degenerate by the same fixture that serves the pixel clause
      ($f_x \ne f_y$, off-axis principal point, nonzero mount station,
      non-identity mount rotation), so no echoed constant is recoverable
      from another. Measured worst residual, \num{1.705303e-13} pixels
      against the $10^{-6}$ gate. \\
    \testid{test_camera_echo_mutation_is_detected} &
      Non-vacuity of the row above, by construction rather than by
      inspection: each echoed constant is perturbed in turn and the
      reconstruction is required to leave the gate tolerance, so a gate
      that read the intrinsics from anywhere but the header could not
      pass. Measured worst residuals --- $f_x$ and $f_y$ scaled by
      $1 + 10^{-6}$: \num{4.75e-4} and \num{5.11e-5} pixels; $c_x$, $c_y$
      shifted \num{1e-5} pixels: \num{1.0e-5} each; the mount quaternion
      perturbed \num{1e-9} in one component: \num{1.55e-6}; the mount
      station shifted \qty{1}{\milli\metre}: \num{1.29e-6}. All ten
      mutations tried, including dropping the mount rotation and the
      mount station entirely, are rejected. \\
\end{longtable}

Exit criterion~7's pose bit-exactness clause is satisfied \emph{by
construction} rather than by a test tolerance: the emitted pose channels
are assignments of the same \texttt{double} values the truth writer
receives, never recomputed (implementation note 2), so no arithmetic
exists that could differ. A reader should scope what that buys. Because
both sides of \testid{test_camera_pose_channels_are_bit_exact_truth} are
read from one log written by one run, the assertion compares a value
against a copy of itself, and it cannot fail on any arithmetic ground. It
retains real power against two specific defects, both live on this
fixture: a mis-indexed truth row, and a hook that logged the
mount-offset-shifted station $\vv{r}_{\mathrm{cam}}$ in place of $\vv{r}$,
which differ here because $\vv{r}^{B}_{\mathrm{cam}} = [0.5, -0.25,
0.125]\,\unit{\metre}$ is nonzero. It is adequate for the clause as the
criterion phrases it, and it proves less than the phrase ``bit-exact''
invites.

The pixel clause of criterion~7 is closed by the two pytest rows above.

\paragraph{Intrinsics clause of exit criterion~7.} The criterion requires
the camera hook's pose \emph{and intrinsics} to equal the \texttt{truth}
channels bit-exactly. Read literally the intrinsics half is ill-posed:
\texttt{truth} carries $\vv{r}$, $\vv{v}$, $\quat{q}_{i2b}$,
$\vv{\omega}^{B}$, and mass, so there is no truth channel an intrinsic
could equal, and $f_x$ is not a quantity the propagator produces. What
the clause can mean, and what is now gated, is that the intrinsics the
core \emph{used} are recoverable from the log and are the ones the logged
pixels were actually computed with.

\testid{test_camera_intrinsics_echo_reproduces_logged_pixels} closes it.
The gate reads $f_x$, $f_y$, $c_x$, $c_y$, $\quat{q}_{b2c}$, and
$\vv{r}^{B}_{\mathrm{cam}}$ from the \texttt{gnc.camera} header echo and
the pose from the \texttt{truth} channels --- both sides off the log,
neither off the mission file --- composes
equations~\eqref{eq:camera:pose} and~\eqref{eq:camera:K} through the blind
NumPy reference, and requires the result to reproduce the logged
\texttt{px\_uv} to the same $10^{-6}$ pixel tolerance as the pixel clause.
This is not the earlier configuration-against-configuration comparison:
the echoed constants are load-bearing, so a header echo that disagreed
with the constants the hook projected through fails the gate by a wide
margin. The mutation evidence is tabulated with the test.

The following remains outstanding:
\begin{itemize}
  \item \textbf{Aberration consistency across hooks.} Pairing the camera
    Sun vector of equation~\eqref{eq:camera:sunvec} against the sun-sensor
    direction of Chapter~\ref{ch:sensors-optical} on one run, confirming
    both resolve the same apparent direction through the shared aberration
    path. No committed test compares the two.
\end{itemize}

\paragraph{Scope of the implemented hook.} Landmarks are body-fixed points
on the \emph{central} body, which is the surface-relative
optical-navigation case this hook targets and is served exactly by the
central-body-fixed rotation the environment model already composes for the
dynamics. Landmarks on a third body, and with them the multi-body
occlusion test of equation~\eqref{eq:camera:occlusion}, require a
configured body table the sensor layer does not yet receive; the near-side
test of equation~\eqref{eq:camera:nearside} is exact for a spherical
central body and a surface landmark, so self-occlusion is handled without
it. Note also that a landmark placed \emph{at} the body center is
degenerate for that test, whose dot product against
$(\vv{\ell} - \vv{r}_T)$ is then identically zero.

\paragraph{Logged channels versus emitted quantities.} Implementation
note 5 lists the visible-bodies list and the Sun vector among the logged
channels, but the normative \texttt{sensors.camera} record layout carries
$t_s$, $\vv{r}$, $\quat{q}_{i2b}$, and the $2L$ interleaved pixel
coordinates only. The Sun vector of equation~\eqref{eq:camera:sunvec} is
computed through the shared aberration path and exposed to in-process
consumers, and the per-landmark visibility flags likewise; extending the
record to carry them would be a log-format change, which this
implementation deliberately does not make. The intrinsics are a different
case: they are constants rather than per-sample quantities, so they are
not record channels either, but v1.3 echoes them once in the header
(Section~\ref{sec:camera:pose}) rather than leaving them in the resolved
configuration alone. Exit criterion~7 gates the pose, the intrinsics, and
the landmark pixels; the record stream carries the first and the third,
the header echo carries the second, and all three clauses are gated.

\section{References}
\label{sec:camera:references}

The pinhole model, intrinsics matrix, and image-coordinate conventions
follow Hartley and Zisserman~\cite{hartley2004}; body radii and
rotational elements follow Archinal et al.~\cite{archinal2018} through
the frames machinery of Chapter~\ref{ch:frames}; the aberration
correction is specified once in Chapter~\ref{ch:sensors-optical} after
Kaplan~\cite{kaplan2005}. Full bibliographic details appear in the
bibliography.
